\documentclass[11pt,a4paper]{article}

\usepackage[margin=2.4cm]{geometry}
\usepackage{fontspec}
\usepackage[french]{babel}
\usepackage{enumitem}
\usepackage{hyperref}
\usepackage{xcolor}
\usepackage{titlesec}
\usepackage{booktabs}
\usepackage{array}
\usepackage{tabularx}
\usepackage{fancyhdr}

\setmainfont{Segoe UI}
\setsansfont{Segoe UI}
\setmonofont{Consolas}

\definecolor{ink}{HTML}{1F2933}
\definecolor{accent}{HTML}{0F6E56}
\definecolor{rule}{HTML}{D9E2EC}
\definecolor{soft}{HTML}{F4F7F6}

\hypersetup{
  colorlinks=true,
  linkcolor=accent,
  urlcolor=accent,
  pdftitle={Guide de déploiement — Mashrou3App},
  pdfauthor={Équipe Mashrou3App}
}

\pagestyle{fancy}
\fancyhf{}
\fancyhead[L]{\small\color{accent} Mashrou3App}
\fancyhead[R]{\small\color{ink!70} Guide de déploiement}
\fancyfoot[C]{\small\thepage}
\renewcommand{\headrulewidth}{0.4pt}
\renewcommand{\headrule}{\hbox to\headwidth{\color{rule}\leaders\hrule height \headrulewidth\hfill}}

\titleformat{\section}{\Large\bfseries\color{accent}}{\thesection}{0.7em}{}
\titleformat{\subsection}{\large\bfseries\color{ink}}{\thesubsection}{0.6em}{}

\setlist[enumerate]{itemsep=0.35em, topsep=0.4em}
\setlist[itemize]{itemsep=0.25em, topsep=0.35em}

\newcommand{\code}[1]{\texttt{#1}}
\newcommand{\pkg}{\code{com.clubaltruisme.mashrou3app}}

\begin{document}

\begin{center}
  {\LARGE\bfseries\color{accent} Guide de déploiement}\\[0.45em]
  {\Large مهندس حامل لكتاب الله}\\[0.25em]
  {\large Mashrou3App}\\[0.8em]
  {\color{ink!75} Document à l'usage de l'équipe projet}\\[0.25em]
  {\color{ink!75} 24 septembre 2026}
\end{center}

\vspace{0.6em}
\noindent\colorbox{soft}{\parbox{\dimexpr\linewidth-2\fboxsep}{
\textbf{Où en est le projet.} Le serveur est déployé. L'application n'est pas encore publiée sur les stores. La prochaine étape est d'installer l'APK de test sur un téléphone Android, de valider les parcours principaux, puis de publier sur le Play Store.
}}

\tableofcontents
\vspace{1em}

\section{Ce qui est déjà en ligne}

\begin{itemize}
  \item Base Supabase du projet \code{mashrou3App} : migrations \code{0080} et \code{0083} à \code{0087} appliquées.
  \item Correctif de sécurité d'inscription : un nouveau compte reçoit le rôle \code{member}. La requête d'invitation superviseur ne renvoie plus le nom ni le groupe.
  \item Functions déployées : \code{activate-invited-account} et \code{send-push}.
  \item Code poussé sur la branche \code{master} du dépôt GitHub.
  \item APK Android de test (profil EAS \code{preview}) construit. Fichier local :\\
        \code{C:\textbackslash Users\textbackslash pc\textbackslash Desktop\textbackslash Mashrou3App-preview.apk}
\end{itemize}

\begin{center}
\begin{tabularx}{\linewidth}{@{}l X@{}}
\toprule
\textbf{Élément} & \textbf{Valeur} \\
\midrule
Nom affiché & مهندس حامل لكتاب الله \\
Version & 1.0.0 \\
Paquet Android et identifiant iOS & \pkg \\
Compte Expo & \code{oumeyma\_elaammari} \\
Profil de test & \code{preview} (APK, distribution interne) \\
Profil stores & \code{production} (AAB Android, incrément automatique) \\
\bottomrule
\end{tabularx}
\end{center}

\section{Installer la version de test sur un téléphone}

Le fichier \code{.apk} est le paquet à installer sur Android. Il ne s'installe pas comme un logiciel Windows.

\subsection{Copie par câble USB}

\begin{enumerate}
  \item Brancher le téléphone au PC en USB.
  \item Dans la notification USB du téléphone, choisir \textbf{Transfert de fichiers}.
  \item Sur le PC, ouvrir le téléphone dans l'explorateur de fichiers.
  \item Copier \code{Mashrou3App-preview.apk} dans le dossier \textbf{Téléchargements} du téléphone.
  \item Sur le téléphone, ouvrir \textbf{Fichiers}, puis \textbf{Téléchargements}, puis l'APK.
  \item Si Android le demande, autoriser l'installation depuis cette source.
\end{enumerate}

\subsection{Téléchargement sur le même Wi-Fi}

Le PC et le téléphone doivent être sur le même réseau. Tant que le PC qui sert le fichier reste allumé, ouvrir dans le navigateur du téléphone :

\begin{center}
\code{http://192.168.0.194:8765/}
\end{center}

Le téléchargement de l'APK démarre. Ouvrir ensuite le fichier téléchargé et accepter l'installation.

\section{Vérifier l'application avant les stores}

Faire ce parcours avec de vrais comptes, sur le téléphone :

\begin{enumerate}
  \item Ouverture de l'application, écran de connexion, session conservée après fermeture.
  \item Inscription d'un membre : le compte reste \code{member}, et ne devient pas administrateur.
  \item Invitation d'un superviseur : activation du compte, puis accès à ses séances.
  \item Une séance, une présence, un message.
  \item Une notification push, application ouverte puis application fermée.
  \item Photo de profil, depuis la caméra et depuis la galerie.
  \item Affichage en arabe, de droite à gauche.
\end{enumerate}

Si une étape échoue, corriger avant la publication. Un correctif d'écran peut ensuite partir en mise à jour Expo. Un correctif de base ou de function se déploie côté Supabase, sans nouvel APK si l'application n'a pas changé.

\section{Préparer le Play Store}

Créer un compte Google Play Console, puis l'application avec exactement ce nom de paquet :

\begin{center}
\pkg
\end{center}

À remplir dans la console avant qu'un utilisateur puisse télécharger l'application :

\begin{itemize}
  \item Nom affiché : مهندس حامل لكتاب الله
  \item Icône, captures d'écran téléphone, texte court et texte long
  \item Adresse publique d'une politique de confidentialité
  \item Questionnaire de classification du contenu
  \item Fiche \textbf{Sécurité des données} : compte, photos, notifications
  \item Pays de distribution
\end{itemize}

Le premier envoi se fait souvent à la main dans Play Console. Les suivants peuvent passer par EAS.

\section{Construire la version de production Android}

Le profil \code{production} produit un fichier \code{.aab}, format exigé par le Play Store, et incrémente le numéro de version.

Avant le build, enregistrer ces variables sur EAS, environnement \textbf{production}. Les valeurs sont celles du fichier \code{.env} à la racine du projet. Ce sont l'adresse publique du projet Supabase et la clé publique \code{anon}.

\begin{itemize}
  \item \code{EXPO\_PUBLIC\_SUPABASE\_URL}
  \item \code{EXPO\_PUBLIC\_SUPABASE\_ANON\_KEY}
\end{itemize}

À la racine du dépôt :

\begin{verbatim}
npx eas-cli build --platform android --profile production
\end{verbatim}

À la fin, la page Expo fournit le fichier \code{.aab}.

\section{Envoyer l'application sur le Play Store}

\subsection{Premier envoi}

\begin{enumerate}
  \item Dans Play Console, ouvrir la piste \textbf{Tests internes}.
  \item Téléverser le fichier \code{.aab}.
  \item Ajouter les adresses e-mail des testeurs.
  \item Publier la piste.
  \item Installer depuis le lien Play, et non depuis l'APK du Bureau.
\end{enumerate}

Quand cette piste est validée, promouvoir la même version vers \textbf{Production} et l'envoyer en examen. L'examen Google prend en général de quelques heures à quelques jours.

\subsection{Envois suivants}

\begin{verbatim}
npx eas-cli submit --platform android --profile production
\end{verbatim}

La première fois, EAS demande une clé de compte de service Google. Ce fichier JSON se crée dans Play Console, section \textbf{Accès API}. Les commandes suivantes réutilisent cette clé.

\section{iPhone}

Cette partie ne concerne que la publication App Store. Elle demande un compte Apple Developer et l'application créée dans App Store Connect avec l'identifiant :

\begin{center}
\pkg
\end{center}

\begin{verbatim}
npx eas-cli build --platform ios --profile production
npx eas-cli submit --platform ios --profile production
\end{verbatim}

EAS gère le certificat et le profil de signature. Apple demande des captures iPhone, la politique de confidentialité et le questionnaire de chiffrement. Le projet déclare déjà que l'application n'utilise pas de chiffrement soumis à déclaration (\code{ITSAppUsesNonExemptEncryption = false}).

\section{Après la publication}

\begin{center}
\begin{tabularx}{\linewidth}{@{}l X@{}}
\toprule
\textbf{Type de changement} & \textbf{Action} \\
\midrule
Texte ou écran, sans module natif &
\code{npx eas-cli update --channel production --message "..."} \\
Icône, permission, notification native, module natif &
Nouveau \code{eas build}, puis \code{eas submit} \\
Base de données ou function &
Migration ou déploiement Supabase. Nouvel binaire seulement si l'application a changé. \\
\bottomrule
\end{tabularx}
\end{center}

\section{Ordre de travail pour l'équipe}

\begin{enumerate}
  \item Copier l'APK de test sur un téléphone Android et l'installer.
  \item Exécuter le parcours de vérification de la section 3.
  \item Corriger les défauts trouvés avant toute publication.
  \item Créer la fiche Play Console et la politique de confidentialité.
  \item Enregistrer les variables Supabase sur l'environnement EAS \code{production}.
  \item Lancer le build \code{production} Android et déposer l'AAB en tests internes.
  \item Après validation, promouvoir la version en production.
  \item Traiter l'App Store seulement si une version iPhone est décidée.
\end{enumerate}

\end{document}
